\section{Faithfulness Scoring Is Metric-Fragile}
\label{sec:metric_fragility}

% POLISH FIX 2: Added paired percentile bootstrap intervals for cross-scorer
% and human-calibration correlations, with the bootstrap protocol stated here.
The same generated answers can tell different stories depending on the
faithfulness observer. We therefore rescore the full \(7{,}500\)-row scaled
cell with DeBERTa-v3, a second NLI model, and a local RAGAS-style judge while
holding the generations fixed.

For correlation uncertainty, we use 10,000 paired-index bootstrap resamples
with a fixed seed and percentile 95\% intervals. Each resample preserves the
pairing between the two scores being correlated; human-label correlations
resample the adjudicated ordinal labels and automatic score together.

\begin{table}[H]
\centering
\small
\caption{Identical generations produce different intervention magnitudes under
different faithfulness judges.}
\label{tab:metrics}
\begin{tabular}{lrrrr}
\toprule
Condition & \(n\) & DeBERTa & Second NLI & RAGAS-style judge \\
\midrule
Baseline & 2500 & 0.661 & 0.350 & 0.730 \\
HCPC-v1 & 2500 & 0.650 & 0.318 & 0.590 \\
HCPC-v2 & 2500 & 0.661 & 0.351 & 0.728 \\
\midrule
Baseline--v1 & -- & 0.011 & 0.032 & 0.140 \\
\bottomrule
\end{tabular}
\end{table}

\begin{table}[H]
\centering
\scriptsize
\caption{Metric agreement on the same \(7{,}500\) generations. Brackets give
paired bootstrap 95\% CIs.}
\label{tab:metric_corr}
\begin{tabular}{lll}
\toprule
Metric pair & Pearson \(r\) [95\% CI] & Spearman \(\rho\) [95\% CI] \\
\midrule
DeBERTa / second NLI & 0.259 [0.240, 0.278] & 0.265 [0.245, 0.285] \\
DeBERTa / RAGAS-style judge & 0.182 [0.162, 0.202] & 0.212 [0.191, 0.233] \\
Second NLI / RAGAS-style judge & 0.674 [0.661, 0.686] & 0.651 [0.638, 0.665] \\
\bottomrule
\end{tabular}
\end{table}

This is the central positive result. The DeBERTa-measured effect is
small; the second NLI effect is moderate; the RAGAS-style judge effect is
roughly an order of magnitude larger than DeBERTa's. Because DeBERTa and
the RAGAS-style judge correlate weakly on identical generations, the
disagreement is not a monotone rescaling. The
\S\ref{sec:human_calibration} human calibration shows that none of the
three scorers is strongly aligned with adjudicated support, so the
disagreement cannot be resolved by privileging the scorer with the
largest measured effect: evaluator choice changes the empirical
conclusion, and no single automatic observer in this audit cleanly
adjudicates that choice.

\subsection{Human Calibration of Metric Disagreement}
\label{sec:human_calibration}

To calibrate the scorer disagreement, two raters annotate 99 examples from the
fixed-generation cell. The raw agreement is 0.919 and Cohen's \(\kappa\) is
0.774. Adjudicated labels are 75 supported, 20 partially supported, and 4
unsupported. Spearman correlation with the adjudicated ordinal labels is weak
for DeBERTa, moderate for the second NLI scorer, and moderate for the
RAGAS-style judge.

\begin{table}[H]
\centering
\scriptsize
\caption{Human calibration audit (\(n=99\), two raters, adjudicated labels).
Brackets give paired bootstrap 95\% CIs.}
\label{tab:human_calibration}
\begin{tabular}{llll}
\toprule
Metric & Spearman \(\rho\) [95\% CI] & Pearson \(r\) [95\% CI] & Notes \\
\midrule
DeBERTa-v3 NLI & 0.140 [-0.047, 0.314] & 0.145 [-0.016, 0.296] & Weak alignment \\
Second NLI & 0.380 [0.208, 0.534] & 0.385 [0.225, 0.523] & Moderate alignment \\
RAGAS-style judge & 0.441 [0.249, 0.624] & 0.338 [0.121, 0.542] & Moderate alignment \\
\bottomrule
\end{tabular}
\end{table}

The human calibration does not settle faithfulness evaluation generally. It
does rule out a comfortable interpretation in which any one automatic metric
can be treated as obviously correct. In this sample, the RAGAS-style judge is
closest to human rank order, but only moderately; DeBERTa is poorly aligned.
The pairwise differences in rank alignment between scorers have overlapping
bootstrap intervals, so the relative ordering should be read as suggestive
rather than definitive. The safe conclusion is that metric disagreement must
be reported, and small single-metric effects should not carry broad claims.
